\documentclass[12pt,a4paper]{article}

% ------------------------------------------------
% Paquetes
% ------------------------------------------------
\usepackage[utf8]{inputenc}
\usepackage[T1]{fontenc}
\usepackage[spanish, es-nodecimaldot]{babel}
\usepackage{geometry}
\usepackage{graphicx}
\usepackage{float}
\usepackage{xcolor}
\usepackage{listings}
\usepackage{hyperref}
\usepackage{fancyhdr}
\usepackage{titlesec}

% ------------------------------------------------
% Configuración de página
% ------------------------------------------------
\geometry{
    top=2.5cm,
    bottom=2.5cm,
    left=2.5cm,
    right=2.5cm
}

% ------------------------------------------------
% Configuración de listings
% ------------------------------------------------
\definecolor{codegray}{gray}{0.95}

\lstset{
    backgroundcolor=\color{codegray},
    basicstyle=\ttfamily\small,
    breaklines=true,
    frame=single,
    columns=fullflexible,
    keepspaces=true,
    showstringspaces=false,
    tabsize=2
}

% ------------------------------------------------
% Encabezado y pie de página
% ------------------------------------------------
\pagestyle{fancy}
\fancyhf{}
\lhead{Interacción de redes}
\rhead{Práctica 6}
\cfoot{\thepage}
\setlength{\headheight}{15pt}

% ------------------------------------------------
% Configuración de títulos
% ------------------------------------------------
\titleformat{\section}
{\large\bfseries}
{\thesection.}
{0.5em}
{}

% ------------------------------------------------
% Hipervínculos
% ------------------------------------------------
\hypersetup{
    colorlinks=true,
    linkcolor=black,
    urlcolor=blue,
    citecolor=black
}

% ------------------------------------------------
% Comando para insertar capturas opcionales
% Si el archivo no existe, se muestra un marcador.
% ------------------------------------------------
\newcommand{\captura}[2]{%
    \begin{figure}[H]
        \centering
        \IfFileExists{#1}{%
            \includegraphics[width=0.92\textwidth,keepaspectratio]{#1}
        }{%
            \fbox{%
                \parbox[c][4cm][c]{0.82\textwidth}{%
                    \centering
                    \textbf{Captura pendiente}\\[0.5em]
                    Archivo esperado: \texttt{#1}
                }%
            }%
        }
        \caption{#2}
    \end{figure}
}

% ------------------------------------------------
% Datos del documento
% ------------------------------------------------
\title{
    Práctica 6: \textbf{OSPF y ACL extendidas}\\[0.5em]
    \large Interacción de redes
}

\author{
    Jorge Alberto Suárez Saldaña\\
    \small Matrícula: 355992\\
    \small Profesor: Raymundo Antonio González Grimaldo
}

\date{\today}

% ------------------------------------------------
% Documento
% ------------------------------------------------
\begin{document}

\maketitle

\begin{abstract}
En esta práctica se implementó una infraestructura de red empresarial
en Cisco Packet Tracer, utilizando enrutamiento dinámico mediante OSPF
y listas de control de acceso extendidas (ACL) para regular el tráfico
entre distintas redes. Se establecieron políticas de comunicación para
las redes de usuarios, administración, servidores y una red externa.
Asimismo, se consideraron los servicios HTTP, HTTPS, FTP y SSH
para comprobar cómo las reglas de filtrado permiten o restringen
las comunicaciones de acuerdo con los requisitos de seguridad.
\end{abstract}

\section{Introducción}

En una infraestructura de red empresarial no basta con garantizar que
los dispositivos puedan comunicarse entre sí. También es necesario
controlar qué dispositivos pueden acceder a determinados recursos y
qué servicios tienen permitido utilizar.

El enrutamiento permite que los paquetes lleguen a redes distintas
de aquella en la que se originaron. Sin embargo, el hecho de que exista
una ruta válida no implica que toda comunicación deba estar autorizada.
Para implementar restricciones se utilizan las listas de control de
acceso, conocidas como ACL (\textit{Access Control Lists}).

En esta práctica se construyó una topología con varios routers,
redes locales, servidores y una red externa. Se empleó OSPF
(\textit{Open Shortest Path First}) para intercambiar rutas entre
los routers internos y se utilizaron ACL extendidas para filtrar
tráfico según las direcciones IP de origen y destino, el protocolo
y, cuando corresponde, los puertos TCP.

Las políticas establecidas buscan permitir las comunicaciones
necesarias para el funcionamiento de la empresa y restringir las
que no están autorizadas. Se prestó especial atención a la protección
de los servidores, al acceso remoto mediante SSH y a la separación
entre el tráfico HTTP y HTTPS.

\section{Objetivos}

\subsection{Objetivo general}

Configurar y evaluar una red empresarial en Cisco Packet Tracer
mediante OSPF y ACL extendidas, aplicando políticas de acceso
entre redes de usuarios, administración, servidores y una red externa.

\subsection{Objetivos específicos}

\begin{itemize}
    \item Diseñar una topología con múltiples routers y redes LAN.
    \item Configurar direccionamiento IPv4 en las interfaces de red.
    \item Implementar OSPF para el intercambio de rutas entre routers.
    \item Configurar una ruta predeterminada hacia la red externa.
    \item Habilitar el acceso remoto seguro mediante SSH en los routers.
    \item Configurar ACL extendidas para controlar el tráfico entre redes.
    \item Permitir HTTPS y bloquear HTTP en las comunicaciones
    de la red empresarial hacia la red externa.
    \item Restringir el acceso a los servidores según la red de origen
    y el servicio solicitado.
    \item Permitir el acceso de administración a los servidores.
    \item Establecer reglas para FTP, HTTPS, ICMP y SSH.
    \item Definir pruebas que permitan comprobar las políticas de acceso.
\end{itemize}

\section{Fundamentos teóricos}

\subsection{Enrutamiento dinámico}

El enrutamiento es el proceso mediante el cual un router determina
por dónde debe enviar un paquete para alcanzar una red de destino.
Las tablas de enrutamiento contienen información sobre las redes
conocidas y el siguiente salto o interfaz que debe utilizarse.

En una red con varios routers, las rutas pueden configurarse
manualmente o aprenderse mediante un protocolo de enrutamiento
dinámico. En esta práctica se utilizó OSPF para las redes internas.

\subsection{Protocolo OSPF}

OSPF es un protocolo de enrutamiento de estado de enlace. Los routers
que participan en el mismo dominio intercambian información sobre
la topología y calculan las rutas mediante el algoritmo SPF
(\textit{Shortest Path First}).

Entre sus características principales se encuentran:

\begin{itemize}
    \item Utiliza un identificador de proceso local para su configuración.
    \item Permite organizar routers y redes en áreas.
    \item Utiliza identificadores de router (\textit{Router ID}).
    \item Calcula rutas considerando el costo de los enlaces.
    \item Actualiza las rutas cuando detecta cambios en la topología.
\end{itemize}

En esta práctica se empleó el área 0, correspondiente al área
de backbone de OSPF. La red externa conectada a R4 no se incorporó
al proceso OSPF interno; en su lugar, se configuró una ruta
predeterminada en R3 y se anunció a los routers internos.

\subsection{Listas de control de acceso extendidas}

Una ACL es una secuencia ordenada de reglas que permite o deniega
paquetes. Una ACL extendida puede evaluar diferentes campos,
entre ellos:

\begin{itemize}
    \item Dirección IP de origen.
    \item Dirección IP de destino.
    \item Protocolo de transporte o de red, como TCP, UDP o ICMP.
    \item Puerto de origen o destino para protocolos como TCP y UDP.
\end{itemize}

Las reglas se procesan de arriba hacia abajo. Cuando un paquete
coincide con una regla, se aplica esa acción y no se continúa
evaluando el resto de la lista. Además, existe una denegación
implícita al final de las ACL de Cisco: si ningún permiso coincide,
el paquete se descarta.

Por esta razón, el orden de las reglas es fundamental. Las
excepciones específicas deben colocarse antes de las denegaciones
generales.

\subsection{Dirección y punto de aplicación de una ACL}

Una ACL puede aplicarse en dirección de entrada (\texttt{in})
o de salida (\texttt{out}) en una interfaz.

\begin{itemize}
    \item \textbf{Entrada (\texttt{in}):} evalúa los paquetes cuando
    ingresan al router por la interfaz correspondiente.
    \item \textbf{Salida (\texttt{out}):} evalúa los paquetes que
    van a salir del router por dicha interfaz.
\end{itemize}

La selección de interfaz y dirección determina qué tráfico será
evaluado. Por ejemplo, una ACL aplicada en salida sobre la interfaz
de R2 que conecta con los servidores puede controlar los paquetes
que el router envía hacia esa red.

Una ACL aplicada en una interfaz no se aplica automáticamente
a todas las demás interfaces del router. Por ello, la ubicación
de cada lista debe corresponder al flujo de tráfico que se desea
controlar.

\subsection{Servicios y puertos utilizados}

En la práctica se contemplaron los siguientes protocolos:

\begin{itemize}
    \item \textbf{HTTP:} utiliza normalmente TCP 80 para el acceso
    a páginas web sin cifrado.
    \item \textbf{HTTPS:} utiliza normalmente TCP 443 para la
    comunicación web cifrada.
    \item \textbf{FTP:} utiliza TCP 20 como canal de data y 21 como canal de control.
    El canal de datos depende del modo de transferencia y de
    la implementación del servicio.
    \item \textbf{SSH:} utiliza TCP 22 para la administración
    remota cifrada.
\end{itemize}

En una ACL, una regla como \texttt{permit tcp ... eq 443} permite
tráfico TCP cuyo puerto de destino coincide con 443. Esto no equivale
a permitir todo el tráfico de una aplicación en cualquier dirección:
la regla se evalúa de acuerdo con las direcciones, el protocolo,
el puerto y la interfaz donde se aplica.

\section{Diseño de la topología}

\subsection{Descripción general}

La topología se organizó en las siguientes redes:

\begin{itemize}
    \item \textbf{RedA:} primera red de usuarios empresariales.
    \item \textbf{RedB:} segunda red de usuarios empresariales.
    \item \textbf{Red de servidores:} contiene Serv1, Serv2 y Serv3.
    \item \textbf{Red de administración:} contiene los equipos
    utilizados para tareas administrativas.
    \item \textbf{Red externa:} contiene PCEXT y ServE.
\end{itemize}

Los routers se conectaron mediante enlaces seriales para permitir
la comunicación entre las distintas redes. R1 conecta las redes
de usuarios; R2 proporciona acceso a la red de servidores; R3
conecta la red de administración y funciona como salida hacia
R4; y R4 conecta la infraestructura interna con la red externa.

\captura{topologia.png}
{Topología completa de la infraestructura implementada en Packet Tracer.}

\section{Plan de direccionamiento IPv4}

Se utilizaron subredes /24 para las redes LAN y subredes /30 para
los enlaces seriales punto a punto. Una máscara /24 permite disponer
de 254 direcciones de host utilizables, mientras que una subred /30
dispone de dos direcciones de host utilizables.

\subsection{Redes LAN}

\begin{table}[H]
\centering
\small
\begin{tabular}{|l|l|l|l|}
\hline
\textbf{Red} & \textbf{Dirección} & \textbf{Gateway} & \textbf{Máscara} \\
\hline
RedA & 148.20.0.0/24 & 148.20.0.1 & 255.255.255.0 \\
\hline
RedB & 148.20.1.0/24 & 148.20.1.1 & 255.255.255.0 \\
\hline
Servidores & 148.20.15.0/24 & 148.20.15.1 & 255.255.255.0 \\
\hline
Administración & 148.20.16.0/24 & 148.20.16.1 & 255.255.255.0 \\
\hline
Red externa & 201.30.20.0/24 & 201.30.20.1 & 255.255.255.0 \\
\hline
\end{tabular}
\caption{Direccionamiento de las redes LAN.}
\end{table}

\subsection{Direcciones de los equipos finales}

\begin{table}[H]
\centering
\small
\begin{tabular}{|l|l|l|}
\hline
\textbf{Equipo} & \textbf{Dirección IPv4} & \textbf{Gateway} \\
\hline
PCA & 148.20.0.100 & 148.20.0.1 \\
\hline
PCB & 148.20.0.101 & 148.20.0.1 \\
\hline
PCC & 148.20.1.100 & 148.20.1.1 \\
\hline
PCD & 148.20.1.101 & 148.20.1.1 \\
\hline
Serv1 & 148.20.15.10 & 148.20.15.1 \\
\hline
Serv2 & 148.20.15.11 & 148.20.15.1 \\
\hline
Serv3 & 148.20.15.12 & 148.20.15.1 \\
\hline
PCE & 148.20.16.100 & 148.20.16.1 \\
\hline
PCF & 148.20.16.101 & 148.20.16.1 \\
\hline
PCEXT & 201.30.20.200 & 201.30.20.1 \\
\hline
ServE & 201.30.20.100 & 201.30.20.1 \\
\hline
\end{tabular}
\caption{Direcciones asignadas a los equipos finales.}
\end{table}

\subsection{Enlaces seriales entre routers}

\begin{table}[H]
\centering
\small
\begin{tabular}{|l|l|l|l|}
\hline
\textbf{Enlace} & \textbf{Red /30} & \textbf{Extremo 1} & \textbf{Extremo 2} \\
\hline
R2--R3 & 148.20.20.0/30 & R2: 148.20.20.1 & R3: 148.20.20.2 \\
\hline
R1--R3 & 148.20.20.4/30 & R1: 148.20.20.5 & R3: 148.20.20.6 \\
\hline
R1--R2 & 148.20.20.8/30 & R1: 148.20.20.9 & R2: 148.20.20.10 \\
\hline
R3--R4 & 201.30.21.0/30 & R4: 201.30.21.1 & R3: 201.30.21.2 \\
\hline
\end{tabular}
\caption{Direccionamiento de los enlaces seriales.}
\end{table}

\section{Configuración del enrutamiento OSPF}

\subsection{Identificadores de router}

Se utilizaron los siguientes identificadores para el proceso OSPF:

\begin{itemize}
    \item R1: \texttt{1.1.1.1}.
    \item R2: \texttt{2.2.2.2}.
    \item R3: \texttt{3.3.3.3}.
\end{itemize}

El proceso OSPF se configuró con el identificador local 1 y el
área 0. Los comandos de referencia para anunciar las redes internas
son los siguientes.

\subsection{Configuración de R1}

\begin{lstlisting}
enable
configure terminal
router ospf 1
 router-id 1.1.1.1
 network 148.20.0.0 0.0.0.255 area 0
 network 148.20.1.0 0.0.0.255 area 0
 network 148.20.20.4 0.0.0.3 area 0
 network 148.20.20.8 0.0.0.3 area 0
end
\end{lstlisting}

\subsection{Configuración de R2}

\begin{lstlisting}
enable
configure terminal
router ospf 1
 router-id 2.2.2.2
 network 148.20.15.0 0.0.0.255 area 0
 network 148.20.20.0 0.0.0.3 area 0
 network 148.20.20.8 0.0.0.3 area 0
end
\end{lstlisting}

\subsection{Configuración de R3}

\begin{lstlisting}
enable
configure terminal
router ospf 1
 router-id 3.3.3.3
 network 148.20.16.0 0.0.0.255 area 0
 network 148.20.20.0 0.0.0.3 area 0
 network 148.20.20.4 0.0.0.3 area 0
end
\end{lstlisting}

Las sentencias \texttt{network} utilizan máscaras wildcard. Por
ejemplo, \texttt{0.0.0.255} corresponde a una red /24, mientras
que \texttt{0.0.0.3} corresponde a una red /30.

Es importante destacar que estas sentencias seleccionan las
interfaces que participarán en OSPF según sus direcciones IP.
En la configuración implementada, el enlace entre R3 y R4 se
mantuvo fuera del proceso OSPF.

\subsection{Ruta predeterminada hacia la red externa}

R3 se configuró para enviar hacia R4 el tráfico que no tuviera
una ruta más específica en su tabla de enrutamiento. La ruta
predeterminada se configuró con el siguiente comando:

\begin{lstlisting}
ip route 0.0.0.0 0.0.0.0 201.30.21.1
\end{lstlisting}

Para anunciar esta ruta al dominio OSPF interno se utilizó:

\begin{lstlisting}
router ospf 1
 default-information originate
\end{lstlisting}

En R4 se configuró una ruta de retorno hacia la red empresarial:

\begin{lstlisting}
ip route 148.20.0.0 255.255.0.0 201.30.21.2
\end{lstlisting}

\section{Configuración de SSH}

\subsection{Propósito}

SSH (\textit{Secure Shell}) permite administrar los routers
remotamente mediante una conexión cifrada. En esta práctica se
habilitó el acceso SSH en R1, R2 y R3 para posteriormente controlar
qué equipo externo podía iniciar una sesión.

La configuración de referencia fue la siguiente:

\begin{lstlisting}
configure terminal
ip domain-name empresa.local
username admin privilege 15 secret ClaveSegura123
crypto key generate rsa
ip ssh version 2
line vty 0 4
 login local
 transport input ssh
end
\end{lstlisting}

El comando de generación de claves RSA puede solicitar el tamaño
del módulo. Se debe utilizar un tamaño admitido por la versión
de IOS simulada; cuando está disponible, es preferible utilizar
un módulo de 2048 bits o superior.

La instrucción \texttt{login local} indica que la autenticación
se realiza con la base de usuarios local del router. La instrucción
\texttt{transport input ssh} limita las conexiones VTY entrantes
a SSH, en lugar de permitir Telnet.

\subsection{Pruebas de acceso remoto}

Desde PCEXT se puede probar la conexión con las direcciones
de las interfaces LAN de los routers:

\begin{lstlisting}
ssh -l admin 148.20.0.1
ssh -l admin 148.20.15.1
ssh -l admin 148.20.16.1
\end{lstlisting}

En el entorno configurado, las credenciales locales utilizadas
fueron:

\begin{itemize}
    \item Usuario: \texttt{admin}.
    \item Contraseña: \texttt{ClaveSegura123}.
\end{itemize}

Las credenciales deben tratarse como información sensible y
no reutilizarse en redes reales sin cambiarlas por contraseñas
adecuadas. En un entorno productivo también conviene aplicar
mecanismos de autenticación y control de acceso más robustos.

\captura{ssh.png}
{Prueba de acceso SSH desde PCEXT hacia los routers.}

\section{Servicios habilitados en los servidores}

Se habilitaron los servicios necesarios para probar las políticas
de acceso en los servidores internos.

\begin{table}[H]
\centering
\small
\begin{tabular}{|l|l|l|}
\hline
\textbf{Servidor} & \textbf{Dirección IP} & \textbf{Servicios considerados} \\
\hline
Serv1 & 148.20.15.10 & HTTP, HTTPS y FTP \\
\hline
Serv2 & 148.20.15.11 & HTTP, HTTPS y FTP \\
\hline
Serv3 & 148.20.15.12 & HTTP, HTTPS y FTP \\
\hline
ServE & 201.30.20.100 & HTTP y HTTPS \\
\hline
\end{tabular}
\caption{Servicios habilitados para las pruebas.}
\end{table}

La disponibilidad de un servicio en el servidor es distinta
de su autorización a través de la red. Por ejemplo, aunque HTTP
esté habilitado en ServE, una ACL puede impedir que los equipos
empresariales establezcan conexiones HTTP con ese servidor.

\section{Políticas de seguridad}

Las políticas de la práctica se resumen en la siguiente tabla.

\begin{table}[H]
\centering
\small
\begin{tabular}{|p{3.1cm}|p{8.8cm}|}
\hline
\textbf{Origen} & \textbf{Política de acceso} \\
\hline
Red empresarial & Puede acceder a la red externa mediante HTTPS,
pero no mediante HTTP. \\
\hline
Red empresarial & Puede enviar solicitudes ICMP de tipo echo
(\texttt{ping}) hacia la red externa. \\
\hline
Red externa & No puede iniciar solicitudes \texttt{ping} hacia
la red empresarial. \\
\hline
PCEXT & Puede establecer conexiones SSH autorizadas hacia R1,
R2 y R3. \\
\hline
Otros equipos externos & No deben poder iniciar conexiones SSH
a los routers protegidos. \\
\hline
Red externa & Solo puede acceder a Serv2 y Serv3 mediante HTTPS
y FTP. \\
\hline
RedA y RedB & Pueden acceder a Serv1 y Serv2 mediante HTTPS y FTP. \\
\hline
RedA y RedB & Pueden acceder a Serv3 mediante FTP, pero no mediante
otros protocolos de aplicación no autorizados. \\
\hline
Red de administración & Puede acceder a todos los servidores mediante
los protocolos contemplados en la política. \\
\hline
Redes internas & Se permite la comunicación ICMP entre equipos
empresariales, de acuerdo con el requisito de conectividad interna. \\
\hline
\end{tabular}
\caption{Resumen de las políticas de filtrado.}
\end{table}

\section{Configuración de las ACL extendidas}

\subsection{ACL de acceso a los servidores en R2}

Se utilizó una ACL extendida llamada \texttt{SERVERS-IN} para
controlar el tráfico que sale de R2 hacia la red de servidores.
La lista se aplicó en dirección de salida sobre la interfaz
GigabitEthernet0/0/0, conectada a dicha red.

La siguiente configuración reúne las reglas correspondientes
a las políticas definidas:

\begin{lstlisting}
configure terminal
ip access-list extended SERVERS-IN

! Red de administracion: acceso a todos los servidores
permit ip 148.20.16.0 0.0.0.255 148.20.15.0 0.0.0.255

! RedA: HTTPS y FTP hacia Serv1
permit tcp 148.20.0.0 0.0.0.255 host 148.20.15.10 eq 443
permit tcp 148.20.0.0 0.0.0.255 host 148.20.15.10 eq 21

! RedA: HTTPS y FTP hacia Serv2
permit tcp 148.20.0.0 0.0.0.255 host 148.20.15.11 eq 443
permit tcp 148.20.0.0 0.0.0.255 host 148.20.15.11 eq 21

! RedA: solamente FTP hacia Serv3
permit tcp 148.20.0.0 0.0.0.255 host 148.20.15.12 eq 21

! RedB: HTTPS y FTP hacia Serv1
permit tcp 148.20.1.0 0.0.0.255 host 148.20.15.10 eq 443
permit tcp 148.20.1.0 0.0.0.255 host 148.20.15.10 eq 21

! RedB: HTTPS y FTP hacia Serv2
permit tcp 148.20.1.0 0.0.0.255 host 148.20.15.11 eq 443
permit tcp 148.20.1.0 0.0.0.255 host 148.20.15.11 eq 21

! RedB: solamente FTP hacia Serv3
permit tcp 148.20.1.0 0.0.0.255 host 148.20.15.12 eq 21

! Red externa: HTTPS y FTP hacia Serv2
permit tcp 201.30.20.0 0.0.0.255 host 148.20.15.11 eq 443
permit tcp 201.30.20.0 0.0.0.255 host 148.20.15.11 eq 21

! Red externa: HTTPS y FTP hacia Serv3
permit tcp 201.30.20.0 0.0.0.255 host 148.20.15.12 eq 443
permit tcp 201.30.20.0 0.0.0.255 host 148.20.15.12 eq 21

! Bloqueo del resto del trafico hacia los servidores
deny ip any 148.20.15.0 0.0.0.255
end
\end{lstlisting}

La ACL se aplicó de la siguiente manera:

\begin{lstlisting}
configure terminal
interface GigabitEthernet0/0/0
 ip access-group SERVERS-IN out
end
\end{lstlisting}

La dirección \texttt{out} hace que la lista evalúe los paquetes
que R2 envía por la interfaz hacia los servidores. Los paquetes
de respuesta de los servidores ingresan por esa interfaz y,
por tanto, no se evalúan con esta misma ACL de salida.

La regla de denegación general se coloca al final para impedir
que otros protocolos o destinos de la red de servidores sean
accesibles por coincidencia con una regla más general.

La regla de administración permite todo el tráfico IP entre
la red de administración y la red de servidores. Por ello,
las reglas específicas posteriores para RedA, RedB y la red
externa no amplían los permisos de administración: esa excepción
ya se concede mediante la primera regla.

\subsection{ACL de tráfico externo en R3}

En R3 se configuró una ACL extendida llamada \texttt{EXTERNAL-IN},
aplicada en entrada sobre la interfaz serial que conecta con R4.
Su finalidad es evaluar el tráfico procedente de la red externa
antes de que se dirija hacia la infraestructura empresarial.

La configuración utilizada fue:

\begin{lstlisting}
configure terminal
ip access-list extended EXTERNAL-IN

! SSH desde PCEXT hacia direcciones de la infraestructura
permit tcp host 201.30.20.200 host 148.20.0.1 eq 22
permit tcp host 201.30.20.200 host 148.20.1.1 eq 22
permit tcp host 201.30.20.200 host 148.20.20.9 eq 22
permit tcp host 201.30.20.200 host 148.20.20.5 eq 22
permit tcp host 201.30.20.200 host 148.20.15.1 eq 22
permit tcp host 201.30.20.200 host 148.20.20.10 eq 22
permit tcp host 201.30.20.200 host 148.20.20.1 eq 22
permit tcp host 201.30.20.200 host 148.20.16.1 eq 22
permit tcp host 201.30.20.200 host 148.20.20.6 eq 22
permit tcp host 201.30.20.200 host 148.20.20.2 eq 22
permit tcp host 201.30.20.200 host 201.30.21.2 eq 22

! HTTPS y FTP externos hacia Serv2 y Serv3
permit tcp 201.30.20.0 0.0.0.255 host 148.20.15.11 eq 443
permit tcp 201.30.20.0 0.0.0.255 host 148.20.15.11 eq 21
permit tcp 201.30.20.0 0.0.0.255 host 148.20.15.12 eq 443
permit tcp 201.30.20.0 0.0.0.255 host 148.20.15.12 eq 21

! Respuestas de conexiones TCP iniciadas desde la empresa
permit tcp 201.30.20.0 0.0.0.255 eq 443 148.20.0.0 0.0.255.255 established
permit tcp 201.30.20.0 0.0.0.255 148.20.0.0 0.0.255.255 established

! Permitir respuestas a ping originados en la empresa
permit icmp 201.30.20.0 0.0.0.255 148.20.0.0 0.0.255.255 echo-reply

! Bloquear ping iniciado desde la red externa
deny icmp any 148.20.0.0 0.0.255.255 echo

! Bloquear el resto del trafico dirigido a la empresa
deny ip any 148.20.0.0 0.0.255.255
permit ip any any
end
\end{lstlisting}

La ACL se aplicó en la interfaz serial correspondiente:

\begin{lstlisting}
configure terminal
interface Serial0/1/1
 ip access-group EXTERNAL-IN in
end
\end{lstlisting}

\subsection{ACL de tráfico empresarial hacia la red externa}

Para controlar las conexiones que salen de la empresa hacia
la red externa se utilizó \texttt{EXTERNAL-OUT}. La lista permite
HTTPS, bloquea HTTP y permite solicitudes ICMP de tipo echo.

\begin{lstlisting}
configure terminal
ip access-list extended EXTERNAL-OUT

! Permitir HTTPS desde la empresa hacia la red externa
permit tcp 148.20.0.0 0.0.255.255 201.30.20.0 0.0.0.255 eq 443

! Bloquear HTTP desde la empresa hacia la red externa
deny tcp 148.20.0.0 0.0.255.255 201.30.20.0 0.0.0.255 eq 80

! Permitir ping iniciado desde la empresa
permit icmp 148.20.0.0 0.0.255.255 201.30.20.0 0.0.0.255 echo

! Permitir el resto del trafico no restringido por las reglas anteriores
permit ip any any
end
\end{lstlisting}

La ACL se aplicó en dirección de salida sobre el enlace de R3 hacia
R4:

\begin{lstlisting}
configure terminal
interface Serial0/1/1
 ip access-group EXTERNAL-OUT out
end
\end{lstlisting}

Las dos ACL de R3 tienen funciones diferentes: \texttt{EXTERNAL-IN}
filtra paquetes que llegan desde R4, mientras que
\texttt{EXTERNAL-OUT} filtra paquetes que R3 envía hacia R4.
Aplicarlas en ambas direcciones permite tratar por separado
el tráfico de entrada y el de salida.

\subsection{Consideraciones sobre el orden de las reglas}

El orden de las entradas de una ACL determina el resultado
del filtrado. Una regla general colocada antes de una excepción
puede impedir que esa excepción llegue a evaluarse.

En particular, la ACL \texttt{EXTERNAL-IN} debe revisarse para
asegurar que ningún equipo externo distinto de PCEXT pueda
iniciar conexiones SSH hacia direcciones de administración
de los routers. La lista contiene permisos explícitos para
PCEXT, pero el permiso final \texttt{permit ip any any} también
puede admitir tráfico dirigido a direcciones que no coincidan
con la denegación de la red 148.20.0.0/16.

Para reforzar esa política, después de las excepciones SSH
autorizadas y antes del permiso final se debe incluir:

\begin{lstlisting}
ip access-list extended EXTERNAL-IN
 deny tcp any any eq 22
\end{lstlisting}

Este ajuste debe realizarse conservando el orden de las reglas:
primero los permisos SSH específicos para PCEXT, después la
denegación de los demás intentos SSH y, finalmente, las reglas
restantes que deban permitirse. En IOS, conviene verificar el
orden efectivo con \texttt{show access-lists}; si es necesario,
se puede reconstruir la ACL para garantizar la posición correcta.

\section{Verificación de las ACL}

\subsection{Comprovación}

\captura{access-listR2.png}
{Salida del comando \texttt{show access-lists} en R2.}

\captura{access-listR3.png}
{Salida del comando \texttt{show access-lists} en R3.}

\section{Pruebas HTTP y HTTPS}

Desde una PC de RedA o RedB se debe intentar acceder al servidor
externo ServE mediante HTTP y HTTPS.

HTTP hacia \texttt{http://201.30.20.100} debe ser bloqueado por la política empresarial.
\captura{http.png}
{http es bloqueado}
    
HTTPS hacia \texttt{https://201.30.20.100} debe ser permitido por la política de salida.
\captura{https.png}
{https es permitido}

Para interpretar el resultado, se debe verificar que los servicios
HTTP y HTTPS estén habilitados en ServE. Si Packet Tracer no
permite completar una sesión HTTPS en el navegador simulado,
se deben revisar los servicios disponibles en esa versión y
complementar la prueba con los contadores de la ACL.

\subsection{Pruebas de conectividad}

\captura{pings.png}
{pings corresponden con lo estipulado por las políticas de seguridad}

\section{Conclusiones}

En esta práctica se trabajó con una infraestructura de red
empresarial compuesta por varias redes LAN y enlaces entre
routers. El direccionamiento IPv4 permitió identificar las
redes de usuarios, servidores, administración y la red externa.

OSPF se utilizó para intercambiar rutas entre los routers
internos, mientras que una ruta predeterminada en R3 permitió
dirigir el tráfico externo hacia R4. En el router de salida
también se contempló una ruta de retorno hacia las redes
empresariales.

La configuración de SSH permitió plantear un mecanismo de
administración remota cifrada. A través de las ACL extendidas
se definieron permisos y restricciones por dirección IP,
protocolo y puerto. Estas reglas permiten expresar políticas
específicas para HTTPS, HTTP, FTP, ICMP y SSH.

La práctica demuestra que el enrutamiento y la seguridad cumplen
funciones diferentes. OSPF determina cómo alcanzar una red,
mientras que las ACL establecen qué tráfico se autoriza a
pasar por las interfaces donde están aplicadas.

Finalmente, se identificó la importancia de validar el orden
de las reglas, el sentido de aplicación de cada ACL y la
configuración de los servicios. La revisión de contadores,
tablas de enrutamiento y pruebas desde los equipos finales
permite distinguir los problemas de conectividad de los
bloqueos deliberados.

\end{document}